% fig:taxonomia — Bubble chart: escopo temporal × substrato × política
% Cada célula mostra a contagem de estudos; cor indica política dominante.
% Requer: tikz (já carregado em main.tex)
\begin{tikzpicture}[
  font=\small,
  cell/.style={
    minimum width=2.0cm, minimum height=1.4cm,
    align=center, inner sep=2pt,
  },
]

% --- Cores de política (aliases de cores definidas em main.tex) ---
\colorlet{polLLM}{archBanco}          % self_control_llm
\colorlet{polHeur}{archHibrido}       % heurística
\colorlet{polApre}{archGit}           % aprendida
\colorlet{polMix}{neutGray}           % célula sem maioria (empate)

% --- Eixos ---
% Colunas: substrato (texto, estruturado, híbrido, executável)
% Linhas: escopo (híbrida, procedural, episódica, semântica)

% Cabeçalhos de coluna
\node[cell, font=\small\bfseries] at (0, 0) {};
\node[cell, font=\small\bfseries] at (2.2, 0) {Texto};
\node[cell, font=\small\bfseries] at (4.4, 0) {Estruturado};
\node[cell, font=\small\bfseries] at (6.6, 0) {Híbrido};
\node[cell, font=\small\bfseries] at (8.8, 0) {Executável};

% Cabeçalhos de linha
\node[cell, font=\small\bfseries, anchor=east] at (-0.2, -1.6) {Híbrida};
\node[cell, font=\small\bfseries, anchor=east] at (-0.2, -3.2) {Procedural};
\node[cell, font=\small\bfseries, anchor=east] at (-0.2, -4.8) {Episódica};
\node[cell, font=\small\bfseries, anchor=east] at (-0.2, -6.4) {Semântica};

% --- Labels de eixo ---
\node[rotate=90, font=\footnotesize\bfseries, anchor=south]
  at (-2.8, -4.0) {Escopo temporal};
\node[font=\footnotesize\bfseries] at (5.5, 0.8)
  {Substrato representacional};

% --- Indicadores de célula vazia (ponto sutil) ---
\foreach \cx/\cy in {
  6.6/-3.2, 8.8/-1.6, 8.8/-4.8, 8.8/-6.4,
  2.2/-6.4, 6.6/-4.8, 6.6/-6.4} {
  \node[circle, fill=black!8, minimum size=0.15cm] at (\cx, \cy) {};
}

% --- Bolhas ---
% Tamanho do círculo proporcional a N; cor = política majoritária na célula

% Híbrida × texto = 5 (Bui, Hosain, Vasilopoulos, Wong, Wu)
% Majoritária LLM (4/5): bui, hosain, wong, wu; heurística: vasilopoulos
\node[circle, fill=polLLM!50, draw=polLLM, minimum size=0.9cm]
  at (2.2, -1.6) {\textbf{5}};

% Híbrida × estruturado = 4 (Chen, Deng, Dong, Zhou)
% Majoritária LLM (3/4): chen, deng, zhou; aprendida: dong
\node[circle, fill=polLLM!50, draw=polLLM, minimum size=0.8cm]
  at (4.4, -1.6) {\textbf{4}};

% Híbrida × híbrido = 5
% (Nadafian, Shen TALM, Tablan, Wang RepoMem, Wang CodeMEM)
% Majoritária LLM (3/5): nadafian, tablan, wang2026codemem;
% heurística: shen2025talm, wang2025repositorymemory
\node[circle, fill=polLLM!50, draw=polLLM, minimum size=0.9cm]
  at (6.6, -1.6) {\textbf{5}};

% Procedural × texto = 1 (Qian) — heurística
\node[circle, fill=polHeur!50, draw=polHeur, minimum size=0.45cm]
  at (2.2, -3.2) {\textbf{1}};

% Procedural × estruturado = 1 (Shen Struct.) — heurística
\node[circle, fill=polHeur!50, draw=polHeur, minimum size=0.45cm]
  at (4.4, -3.2) {\textbf{1}};

% Procedural × executável = 3 (Li GBT, Xia Live-SWE, Zhang DGM)
% Empate: 1 heurística (li), 1 LLM (xia), 1 aprendida (zhang DGM)
\node[circle, fill=polMix!40, draw=polMix, minimum size=0.7cm]
  at (8.8, -3.2) {\textbf{3}};

% Episódica × texto = 3 (Lindenbauer, Su, Zhang AccelOpt) — heurística (3/3)
\node[circle, fill=polHeur!50, draw=polHeur, minimum size=0.7cm]
  at (2.2, -4.8) {\textbf{3}};

% Episódica × estruturado = 1 (Guo EET) — LLM
\node[circle, fill=polLLM!50, draw=polLLM, minimum size=0.45cm]
  at (4.4, -4.8) {\textbf{1}};

% Semântica × estruturado = 1 (Wang MemGovern) — LLM
\node[circle, fill=polLLM!50, draw=polLLM, minimum size=0.45cm]
  at (4.4, -6.4) {\textbf{1}};

% --- Legenda ---
\node[anchor=west] at (0.0, -7.8) {
  \begin{tikzpicture}[font=\footnotesize, baseline=-0.5ex]
    \node[circle, fill=polLLM!50, draw=polLLM,
      minimum size=0.35cm] at (0,0) {};
    \node[anchor=west] at (0.3,0) {LLM controlador};
    \node[circle, fill=polHeur!50, draw=polHeur,
      minimum size=0.35cm] at (3.3,0) {};
    \node[anchor=west] at (3.6,0) {Heurística};
    \node[circle, fill=polMix!40, draw=polMix,
      minimum size=0.35cm] at (5.8,0) {};
    \node[anchor=west] at (6.1,0) {Sem maioria};
  \end{tikzpicture}
};

\end{tikzpicture}
